%-------------------------
% Resume in Latex
% Rasool Peykarporsan
%------------------------

\documentclass[letterpaper,11pt]{article}

\usepackage{latexsym}
\usepackage[empty]{fullpage}
\usepackage{titlesec}
\usepackage{marvosym}
\usepackage[usenames,dvipsnames]{color}
\usepackage{verbatim}
\usepackage{enumitem}
\usepackage[hidelinks]{hyperref}
\usepackage{fancyhdr}
\usepackage[english]{babel}
\usepackage{tabularx}
\usepackage{fontawesome5}
\usepackage[scale=0.90,lf]{FiraMono}

\definecolor{light-grey}{gray}{0.83}
\definecolor{dark-grey}{gray}{0.3}
\definecolor{text-grey}{gray}{.08}

\DeclareRobustCommand{\ebseries}{\fontseries{eb}\selectfont}
\DeclareTextFontCommand{\texteb}{\ebseries}

\usepackage{contour}
\usepackage[normalem]{ulem}
\renewcommand{\ULdepth}{1.8pt}
\contourlength{0.8pt}
\newcommand{\myuline}[1]{%
  \uline{\phantom{#1}}%
  \llap{\contour{white}{#1}}%
}

\usepackage{tgheros}
\renewcommand*\familydefault{\sfdefault}
\usepackage[T1]{fontenc}

\pagestyle{fancy}
\fancyhf{}
\fancyfoot{}
\renewcommand{\headrulewidth}{0pt}
\renewcommand{\footrulewidth}{0pt}

\addtolength{\oddsidemargin}{-0.5in}
\addtolength{\evensidemargin}{0in}
\addtolength{\textwidth}{1in}
\addtolength{\topmargin}{-.5in}
\addtolength{\textheight}{1.0in}

\urlstyle{same}
\raggedbottom
\raggedright
\setlength{\tabcolsep}{0in}

\titleformat{\section}{
    \bfseries \vspace{2pt} \raggedright \large
}{}{0em}{}[\color{light-grey}{\titlerule[2pt]} \vspace{-4pt}]

%-------------------------
% Custom commands
\newcommand{\resumeItem}[1]{
  \item\small{{#1 \vspace{-1pt}}}
}

\newcommand{\resumeSubheading}[4]{
  \vspace{-1pt}\item
    \begin{tabular*}{\textwidth}[t]{l@{\extracolsep{\fill}}r}
      \textbf{#1} & {\color{dark-grey}\small #2}\vspace{1pt}\\
      \textit{#3} & {\color{dark-grey}\small #4}\\
    \end{tabular*}\vspace{-4pt}
}

\newcommand{\resumeProjectHeading}[2]{
    \item
    \begin{tabular*}{\textwidth}{l@{\extracolsep{\fill}}r}
      #1 & {\color{dark-grey}\small #2} \\
    \end{tabular*}\vspace{-4pt}
}

\newcommand{\resumeSubItem}[1]{\resumeItem{#1}\vspace{-4pt}}

\renewcommand\labelitemii{$\vcenter{\hbox{\tiny$\bullet$}}$}

\newcommand{\resumeSubHeadingListStart}{\begin{itemize}[leftmargin=0in, label={}]}
\newcommand{\resumeSubHeadingListEnd}{\end{itemize}}
\newcommand{\resumeItemListStart}{\begin{itemize}}
\newcommand{\resumeItemListEnd}{\end{itemize}\vspace{0pt}}

\color{text-grey}

%-------------------------------------------
\begin{document}

%----------HEADING----------
\begin{center}
    \textbf{\Huge Rasool Peykarporsan} \\ \vspace{5pt}
    \small
    \faGlobe \hspace{2pt} \href{https://rasoolpeykarporsan.me/}{\texttt{rasoolpeykarporsan.me}} \hspace{1pt} $|$
    \hspace{1pt} \faPhone* \hspace{1pt} \texttt{+64 29 023 55104} \hspace{1pt} $|$
    \hspace{1pt} \faEnvelope \hspace{2pt} \href{mailto:Rasool.peykarporsan@gmail.com}{\texttt{Rasool.peykarporsan@gmail.com}} \\
    \faGithub \hspace{2pt} \href{https://github.com/Rasoolpey}{\texttt{Rasoolpey}} \hspace{1pt} $|$
    \hspace{1pt} \faGraduationCap \hspace{2pt} \href{https://scholar.google.com/citations?user=AWOKpCAAAAAJ}{\texttt{Google Scholar}} \hspace{1pt} $|$
    \hspace{1pt} \faMapMarker* \hspace{2pt}\texttt{Auckland, New Zealand}
    \\ \vspace{-3pt}
\end{center}

%-----------ABOUT ME-----------
\section{ABOUT ME}
My journey through a B.Sc. and M.Sc. in Electrical Engineering gave me a solid foundation in analog and digital circuit design, power electronics, control theory, PCB layout, and embedded C/C++. 
Pursuing a PhD at Auckland University of Technology then further strengthened that foundation: 
I learned to see every design from multiple angles, question assumptions, and build reliable solutions that actually stand up to rigorous testing. 
That process also shaped the engineer I am today, someone who picks up new topics quickly, stays self-motivated through long and uncertain projects, and genuinely enjoys working in a collaborative team where every member can grow and contribute. 
I am currently seeking an internship where I can apply my technical knowledge and hands-on, problem-solving mindset to advanced wireless power and circuit design challenges at Apple.

%-----------TECHNICAL SKILLS-----------
\section{TECHNICAL SKILLS}
\begin{itemize}[leftmargin=0in, label={}]
  \small{\item{
    \textbf{Power Electronics \& Hardware}{: Direct wireless charging, resonant \& soft-switching converters (LLC, Class-D), analog \& digital circuit design, PCB layout (Altium, OrCAD)}\vspace{2pt} \\
    \textbf{Circuit Simulation \& Modelling}{: Cadence, ADS, LTspice, MATLAB/Simulink, circuit-level simulation, system-level impedance \& loop-gain shaping}\vspace{2pt} \\
    \textbf{Instrumentation \& Testing}{: Impedance analyser, Virtual Network Analyser (VNA), oscilloscopes, OPAL-RT, dSPACE, STM32F4/F7, Texas Instruments C2000, failure analysis, root cause \& corrective action, DOE creation \& execution}\vspace{2pt} \\
    \textbf{Prototyping \& Validation}{: Hands-on experience (soldering, circuit measurement, hardware prototyping), system bring-up, characterisation, factory data collection support}\vspace{2pt} \\
    \textbf{Programming \& Scripting}{: Python, MATLAB, C, C++, embedded programming, Verilog/HDL}\vspace{2pt} \\
    \textbf{Data \& Statistical Analysis}{: Statistical tools (MATLAB), validation data analysis, PyTorch, Reinforcement Learning, ML-driven fault classification}
  }}
\end{itemize}

%-----------EDUCATION-----------
\section{EDUCATION}
  \resumeSubHeadingListStart

    \resumeSubheading
      {Auckland University of Technology (AUT)}{2023 -- present}
      {PhD --- Power Electronics \& Power Systems}{Auckland, New Zealand}

    \resumeSubheading
      {Shahid Beheshti University (top 5 in Iran)}{2016 -- 2019}
      {M.Sc. --- Power Electronics \& Electric Machines Engineering}{Tehran, Iran}

    \resumeSubheading
      {Ferdowsi University of Mashhad (top 10 in Iran)}{2011 -- 2015}
      {B.Sc. --- Electrical Power Engineering}{Mashhad, Iran}

  \resumeSubHeadingListEnd


%-----------EXPERIENCE-----------
\section{EXPERIENCE}
  \resumeSubHeadingListStart

    \resumeSubheading
      {Auckland University of Technology}{2023 -- present}
      {PhD Researcher \& Teaching Assistant}{Auckland, New Zealand}
      \resumeItemListStart
        \resumeItem{Teaching Assistant across undergraduate and postgraduate courses: Power Systems, Power Electronics, Renewable Energy Resources.}
        \resumeItem{\textbf{Supervised 18 students}: 5 final-year theses and 13 semester research projects.}
        \resumeItem{Peer Mentor for bachelor's, master's, and PhD students (AUT Peer Mentoring Programme).}
      \resumeItemListEnd

    \resumeSubheading
      {Circuit Intelligent}{2025 -- present}
      {Founder}{Auckland, New Zealand}
      \resumeItemListStart
        \resumeItem{Founded an AI-assisted circuit-design platform pairing classical EDA workflows with LLM-driven guidance to shorten the path from design intent to a working schematic.}
      \resumeItemListEnd

    \resumeSubheading
      {Khorasan Steel Complex}{2022 -- 2023}
      {Electrical Engineer, R\&D}{Iran}
      \resumeItemListStart
        \resumeItem{R\&D on Substation design and high-voltage infrastructure for plant expansion, and cathodic protection for pipeline and structural assets.}
      \resumeItemListEnd

  \resumeSubHeadingListEnd


%-----------PUBLICATIONS (SELECTED)----------
\section{SELECTED PUBLICATIONS}
  \resumeSubHeadingListStart

    \resumeProjectHeading
      {\textbf{Intelligent Fractional-Order Cascade Control for Wind Farms (DFIG)} $|$ \textit{IEEE Trans.\ Energy Convers.}}{2025}
      \resumeItemListStart
        \resumeItem{4DoF fractional-order cascade controller tuned by DDPG reinforcement learning; validated on a NZ wind farm.}
      \resumeItemListEnd

    \resumeProjectHeading
      {\textbf{Low-Voltage Solid-State DC Circuit Breaker (BiTriCap)} $|$ \textit{IEEE Trans.\ Power Electron.}}{2024}
      \resumeItemListStart
        \resumeItem{Novel bidirectional thyristor-capacitor suppressor; experimentally validated at 48 VDC / 8 A with only 2 V residual overvoltage.}
      \resumeItemListEnd

    \resumeProjectHeading
      {\textbf{Bidirectional Passive-Technique SSCB for LV Applications} $|$ \textit{IEEE JESTPE}}{2025}
      \resumeItemListStart
        \resumeItem{RCD-based passive SSCB: 20 \textmu s full interrupt, 90 \textmu s total fault clearing at 48 VDC / 10 A.}
      \resumeItemListEnd

    \resumeProjectHeading
      {\textbf{Hierarchical Transformer Fault Diagnosis for Four-Leg Inverters} $|$ \textit{IJEPES}}{2026}
      \resumeItemListStart
        \resumeItem{Two-level Transformer-based open/short-circuit switch fault classifier; dataset generated on OPAL-RT HIL.}
      \resumeItemListEnd

    \resumeProjectHeading
      {\textbf{Data-Driven Port-Hamiltonian Parameter ID of Synchronous Generators} $|$ \textit{Sensors}}{2026}
      \resumeItemListStart
        \resumeItem{Three-stage framework with differential damping decoupling; 8 parameters identified with 1.26--9.10\% error.}
      \resumeItemListEnd

    \resumeProjectHeading
      {\textbf{Model-Free Cyber-Attack Defence for Load Frequency Control} $|$ \textit{IET GTD}}{2026}
      \resumeItemListStart
        \resumeItem{DRL-tuned model-free observer detects and mitigates FDI attacks; validated in real time on OPAL-RT.}
      \resumeItemListEnd

  \resumeSubHeadingListEnd
  \vspace{2pt}
  \small{Full list: \href{https://scholar.google.com/citations?user=AWOKpCAAAAAJ}{\myuline{Google Scholar}} \hspace{3pt} | \hspace{3pt} ORCID: \href{https://orcid.org/0009-0006-5237-3847}{\myuline{0009-0006-5237-3847}}}


%-----------PROJECTS-----------
\section{KEY PROJECTS}
    \resumeSubHeadingListStart

      \resumeProjectHeading
        {\textbf{PHPS --- Port-Hamiltonian Power System Simulator} $|$ \textit{Python, C++}}{github.com/Rasoolpey/PHPS}
        \resumeItemListStart
          \resumeItem{Transient stability simulator with built-in Lyapunov function at every simulation step; Python interface over a compiled C++ solver.}
        \resumeItemListEnd

      \resumeProjectHeading
        {\textbf{Microgrid Hardware Testbed} $|$ \textit{Altium Designer, PCB}}{github.com/Rasoolpey/Microgrid\_hardware\_description}
        \resumeItemListStart
          \resumeItem{Modular converter family (Buck, VSC, NPC, four-leg) designed end-to-end; sensor-rich boards with I2C/SPI ADCs serving as the lab's shared hardware platform.}
        \resumeItemListEnd

      \resumeProjectHeading
        {\textbf{TCP PF--MATLAB--Python Bridge} $|$ \textit{Python, MATLAB, PF}}{github.com/Rasoolpey/TCP-PowerFactory-Matlab-Python}
        \resumeItemListStart
          \resumeItem{Co-simulation framework synchronising DIgSILENT PowerFactory and MATLAB/Simulink with any TCP-capable client; enables ML and RL to drive live simulations.}
        \resumeItemListEnd

      \resumeProjectHeading
        {\textbf{PIGNN --- Physics-Informed GNN Digital Twin} $|$ \textit{Python, PyTorch}}{github.com/Rasoolpey/PIGNN}
        \resumeItemListStart
          \resumeItem{GNN that encodes grid topology and physics equations; learnable components adapt to wear and operational changes.}
        \resumeItemListEnd

    \resumeSubHeadingListEnd


%-----------HONORS AND AWARDS-----------
\section{HONORS \& AWARDS}
  \resumeSubHeadingListStart

    \resumeProjectHeading
      {\textbf{MBIE FAN Project Doctoral Scholarship} $|$ \textit{Ministry of Business, Innovation and Employment (NZ)}}{2021}
      \resumeItemListStart
        \resumeItem{Awarded a fully-funded PhD scholarship as part of the MBIE FAN (Future Architecture of the Network) project.}
      \resumeItemListEnd

  \resumeSubHeadingListEnd


%-----------RECOMMENDATIONS-----------
\section{REFERENCES}
 \begin{itemize}[leftmargin=0in, label={}]
    \small{\item{
     \textbf{Prof. Tek Tjing Lie} $|$ \textit{PhD Supervisor, Head of School at the Auckland University of Technology School of Engineering, Computer and Mathematical Sciences} \\
     \href{tek.lie@aut.ac.nz}{\myuline{tek.lie@aut.ac.nz}} \hspace{2pt}  $|$ \hspace{2pt} \texttt{+64 21813661} \vspace{4pt} \\
     \textbf{Dr. Zoey Zhou} $|$ \textit{Lecturer, Auckland University of Technology} \\
     \href{zoey.zhou@aut.ac.nz}{\myuline{zoey.zhou@aut.ac.nz}} \hspace{2pt}  $|$ \hspace{2pt} \texttt{+64 AUText-31776}
    }}
 \end{itemize}

\end{document}
